\documentclass[11pt]{article}
\usepackage[margin=1in]{geometry}
\usepackage{amsmath,amssymb,amsfonts,mathtools}
\usepackage{array,booktabs,enumitem,graphicx,xcolor,hyperref}
\usepackage[T1]{fontenc}
\usepackage[utf8]{inputenc}
\title{Ashtekar's Approach to Quantum Gravity}
\author{Source-backed arXiv sample}
\date{}
\begin{document}
\maketitle
\begin{abstract}
A review is given of work by Abhay Ashtekar and his colleagues Carlo Rovelli, Lee Smolin, and others, which is directed at constructing a nonperturbative quantum theory of general relativity.
\end{abstract}
\section{Source Metadata}
This sample is based on arXiv record \texttt{hep-th/9109002}. Authors: Gary T. Horowitz.
\section{Extracted Prose 1}
Although the results that have been obtained so far are promising, there is much that remains to be done before one can claim to have a consistent quantum theory of gravity. This section is divided into three parts. In the first, I consider some open questions in the main program described above. The second includes a short discussion of other approaches to quantum gravity using Ashtekar variables. In the third I consider the question of whether there exists yet another set of canonical variables for general relativity (or a theory containing general relativity) such that the constraints are simplified even further.
\section{Extracted Prose 2}
I wish to thank the organizers of the Strings and Symmetries 1991 Conference for a stimulating meeting. I am grateful to A. Ashtekar, C. Rovelli, and L. Smolin for discussions of their work. I have also benefited from comments and conversations with S. Giddings, J. Hartle, T. Jacobson, M. Srednicki, and A. Strominger. This work was supported in part by NSF Grant PHY-9008502.
\end{document}
